\section{Conclusion}
\label{conclusion}

We have devised a POMDP model of penetration testing that allows to
naturally represent many of the features of this application, in
particular incomplete knowledge about the network configuration, as
well as dependencies between different attack possibilities, and
firewalls. Unlike any previous methods, the approach is able to
intelligently mix scans with exploits. While this accurate solution
does not scale, large networks can be tackled by a decomposition
algorithm. Our present empirical results suggest that this is
accomplished at a small loss in quality relative to a global POMDP
solution.

An important open question is to what extent our POMDP + decomposition
approach is more cost-effective than the classical planning solution
currently employed by Core Security. Our next step will be to answer
this question experimentally, comparing the attack quality of 4AL
against that of the policy that runs extensive scans and then attaches
FF's plan for the most probable configuration.

4AL is a domain-specific algorithm and, as such, does not contribute
to the solution of POMDPs in general. At a high level of abstraction,
its idea can be understood as imposing a template on the policy
constructed, thus restricting the space of policies explored (and
employing special-purpose algorithms within each part of the
template). In this, the approach is somewhat similar to known POMDP
decomposition approaches (e.g.,
\cite{PinGorThr-uai03,mueller:biundo:ki-11}). It remains to be seen
whether this connection can turn out fruitful for either future work
on attack planning, or POMDP solving more generally.

The main directions for future work are to devise more accurate models
of software updates (hence obtaining more realistic designs of the
initial belief); to tailor POMDP solvers to this particular kind of
problem, which has certain special features, in particular the absence
of non-deterministic actions and that some of the uncertain parts of
the state (e.g.\ the operating systems) are static; and to drive the
industrial application of this technology. We hope that these will
inspire other researchers as well.


